% !PRINTED 78
\weberpage{78}

All primary divisors of the full cyclotomic field $\Omega_m$ are accordingly representable in the form $\Omega(\eta)$, i.\,e. all numbers of such a divisor are rationally representable by the cyclotomic periods $\eta$, and since one finds the non-primary divisors of $\Omega_m$ again as primary divisors of lower cyclotomic fields, it follows:

\gesperrt{4. All cyclotomic fields are representable in the form $\Omega(\eta)$.}

Or also:

\gesperrt{5. Every rational function of roots of unity can be represented as a rational function of a cyclotomic period.}

And this theorem can also be expressed in the form:

\gesperrt{6. If $\mathfrak{A}$ is an arbitrary primary or non-primary divisor of $\mathfrak{N}$ of index $e$, then there is a divisor $m_1$ of $m$ and a cyclotomic period $\eta$ of index $e$ lying in $\Omega_{m_1}$, such that $\eta$, regarded as a function of $r$, is a function belonging to the group $\mathfrak{A}$, and thus, if $a$ belongs to $\mathfrak{A}$, remains unchanged under the substitution $(r, r^a)$ and receives $e$ different values for all the substitutions of $\mathfrak{N}$.}

If $a$ runs through a group $\mathfrak{A}$ which is a primary divisor of $\mathfrak{N}$, and we set

\[
\eta_h = \sum^a r^{ah},
\]

even when $h$ is not relatively prime to $m$, then we can represent these quantities $\eta_h$, which all admit the permutations of the group $\mathfrak{A}$, rationally by $\eta = \sum r^a$. On the other hand, however, every rational function of $\eta$ can also be represented linearly by the $\eta_h$. This will be proved once it is shown that the product of two $\eta_h$ can be expressed linearly by the $\eta_h$. But we have:

\[
\eta_h \eta_k = \sum^{a, a'} r^{ha + ka'};
\]

if one first takes the sum over $a'$ with $a$ held fixed, one can replace $a'$ by $aa'$, since $aa'$ runs through the group $\mathfrak{A}$ at the same time as $a'$. One thus obtains

\[
\eta_h \eta_k = \sum^{a, a'} r^{a(h + ka')} = \sum^{a'} \eta_{h + ka'}.
\]
